\section{Discussion}
\label{sec:discussion}

\paragraph{What the estimate means.} The singleton-free estimate does not count ``all features a model has''. Capture--recapture cannot answer that, for animals or for features \citep{link2003nonidentifiability}. It counts the features that SAEs of a given size find with probability at least $1/K$ per seed. Two properties make it useful. First, in every synthetic setting where the answer is known and reachable, it is right to within a few percent while all naive counts are badly wrong (Table~\ref{tab:synth}). Second, its behaviour across dictionary sizes is diagnostic. It rises while the dictionary is too small and levels off at the feature count once the dictionary is large enough. A practitioner can therefore read off when a larger SAE stops buying new features, which is the question behind choosing a dictionary size, without any ground truth.

\paragraph{Why seed disagreement is informative.} \citet{paulo2025sparse} showed that seeds disagree. Our results separate two sources of disagreement. Some features are real but rare, so a dictionary too small to hold them all catches different subsets on different seeds. Others are junk directions that a single seed learns. The first kind is structure the census should count. The second is noise it should ignore. The capture-frequency histogram tells them apart, because junk is almost never reproduced.

\paragraph{Limitations.}
\begin{itemize}\setlength\itemsep{1pt}
\item \emph{Synthetic world.} Our main planted features are independent, equally sized random directions with Zipf-like frequencies. Real features are correlated, hierarchical and split as dictionaries grow \citep{bricken2023monosemanticity,leask2025sparse}, which can turn one ``species'' into several. Section~\ref{sec:hier} tests one simple form of splitting, a single parent--child level with fixed geometry. There the census levels off at the finest level only once it fits, and undercounts by \HierSfErrFull{} when the dictionary is just large enough. Real hierarchies are deeper and less regular. The synthetic results show the estimator works when a fixed population exists. They cannot show that one exists in a real model.
\item \emph{Seeds are not independent nets.} All seeds see the same data in the same order, and identification relies on a cosine threshold. We checked $\tau\in\{\TauLo,\MainTau,\TauHi\}$, but correlated failures, such as all seeds merging the same two features, would be invisible to any capture--recapture estimator.
\item \emph{The probability floor defines the target.} With $K$ seeds, features caught with probability well below $1/K$ are out of reach by construction. Halving the floor changed results little in our synthetic settings, but that need not hold elsewhere.
\item \emph{Scale and training budget.} We used dictionaries of at most 8k latents on one layer of a 70M-parameter model, with \RealSeeds{} seeds, trained on \RealTrainTokens{} tokens and checked once with \RealLongTokens{}, on a laptop CPU. More data raised the count, so our Pythia numbers describe SAEs trained on this budget. Frontier SAEs have millions of latents. The method's cost is $K$ training runs, which is significant at that scale.
\item \emph{Few-seed bias.} Fitted to four seeds, the estimator over-predicts what further seeds find by up to about a tenth (\S\ref{sec:heldout}).
\end{itemize}

\paragraph{Broader impact.} This is a measurement tool for interpretability research. Knowing how many features a model has, or that a dictionary has stopped finding new ones, helps decide how much of a model an interpretability method has actually covered. That matters when SAE features are used to audit model behaviour. A risk is false reassurance: a plateau means an SAE family stopped finding \emph{reachable} features, not that the model has no others. We see no direct path to misuse.

